% Generated by build/make_tables.py from exceptions_vs_disclosed.csv. Do not edit.

\begin{table}[tbp]
\centering
\small
\caption{Exception counts at 95\%: the proxy book against Goldman's disclosed counts}
\label{tab:exceptions}
\begin{tabular}{@{}lrrrr@{}}
\toprule
 & \multicolumn{3}{c}{Proxy book, plain historical simulation} & Goldman \\
\cmidrule(lr){2-4}
Year & Days & Exceptions & Rate (\%) & Exceptions \\
\midrule
2016 & 233 & 9 & 3.9 & 0 \\
2017 & 232 & 5 & 2.2 & 0 \\
2018 & 230 & 11 & 4.8 & 2 \\
2019 & 225 & 8 & 3.6 & 0 \\
2020 & 227 & 12 & 5.3 & 2 \\
2021 & 231 & 11 & 4.8 & 1 \\
2022 & 230 & 70 & 30.4 & 2 \\
2023 & 233 & 19 & 8.2 & 1 \\
2024 & 228 & 11 & 4.8 & 2 \\
2025 & 224 & 4 & 1.8 & 0 \\
\midrule
Total & 2,293 & 160 & 6.98 & 10 \\
\bottomrule
\end{tabular}
\notetext{\textbf{These two columns do not measure the same thing and should not be compared as though they did.} Goldman scores VaR against daily net revenues, which include intraday market-making income that is positive far more often than not; the proxy book has no client flow and no bid--offer spread, so its profit and loss is a clean mark on a static position. The table is reported for completeness and the gap is an indication of what the revenue cushion is worth, not evidence that either model is miscalibrated. No conclusion in this paper rests on it.}
\end{table}
